\subsection{The literal two-point elliptic and five-point branch marking}

\begin{result}{E.183}{Proved}
{The normalized two-point action has the original branch parameter}
Let \(F=(1,f)\) be a full GT element with \(\rho(F)=0\).
Identify the ordered two-point torus group with
\[
 \Gamma_{1,2}=B_4/\langle\Omega\rangle,\qquad
 B_4=\langle a,b,c\mid aba=bab,\ bcb=cbc,\ ac=ca\rangle,
 \qquad\Omega=(abc)^4,
\tag{E.183.1}
\]
using the double cover of the four-branch-point disc, with
its two boundary circles capped to the two ordered poles.
Under this actual marking the branch-ordered subgroup is
\(P_4/\langle\Omega\rangle=\Gamma_{0,5}\), of index \(24\).

The action of E.182 on \(\widehat\Gamma_{1,2}\) agrees,
up to one explicit pure inner conjugation, with the standard
genus-zero GT action on the same \(\widehat B_4/\langle\Omega\rangle\).
In particular its restriction to the actual branch-ordered
cover has parameter exactly \(F\).
This proves the two-point branch-order comparison in the
normalized scope of E.181; it does not yet identify the
action obtained from a higher-point involution normalizer.
\end{result}

\begin{proof}
The group identification in E.183.1 is the ordinary
double-cover marking, also used in
\citet[\S1, p.659]{LitSchneps2006Automorphisms}.
The braids \(a,b,c\) lift to the positive twists in a
three-curve chain on the two-boundary torus.  Capping both
boundary circles gives the ordered two-point surface.
The full disc twist becomes trivial after this cap, and
the resulting group is exactly the indicated quotient.
Ordering all four branch points is the index-\(24\)
pure-braid cover.  All these are ordinary maps before
completion, so their finite-index topology is full.

Put
\[
        \omega=(ab)^3,\qquad d=\omega^2=(ab)^6,
        \qquad y_3=ba^2b=\omega a^{-2}.
\tag{E.183.2}
\]
The curve of \(d\) is the boundary of a regular
neighbourhood of the first two chain curves.
The curve decompositions
\[
        P=\{a,c\}\ \xrightarrow{\mathrm A}\
        Q=\{a,d\}\ \xrightarrow{\mathrm S}\
        R=\{b,d\}
\tag{E.183.3}
\]
are literal pants moves: the first replacement curves meet
twice with algebraic intersection zero, and the second
meet once inside the one-holed torus bounded by \(d\).
No branch labels are permuted by the choices in this path.
The flat action based at \(P\) therefore fixes \(a,c\)
and sends \(b\) to its conjugate by
\[
                   W=f(c,d)^{-1}f(a^2,b^2)^{-1}.
\tag{E.183.4}
\]
This is exactly the A-then-S weight order in E.182.2.

The following braid identity is the essential mixed marking:
\[
                 (\omega c)^2=\Omega a^2,
                 \qquad (c\omega)^2=(\omega c)^2.
\tag{E.183.5}
\]
For completeness a defining-relation certificate for the
first equality is displayed below.  Digits \(1,2,3\)
denote \(a,b,c\); each consecutive pair differs by
\(121=212\), \(232=323\), or \(13=31\).
\[
 \begin{gathered}
 12121231212123=12121231211213=12121232121213\\
 =12121232112113=12121232112131=12121232112311\\
 =12121323112311=12121321312311=12123121312311\\
 =12123212312311=12132312312311=12312312312311.
 \end{gathered}
\tag{E.183.6}
\]
The element \(a\) commutes with both \(\omega\) and \(c\),
and \(\Omega\) is the full braid twist and is central.
Thus the first expression in E.183.5 commutes with
\(\omega\); conjugating it by \(\omega^{-1}\) proves
the second equality.  Also
\(a^2ba^2b=(ab)^3\), by one use of the braid relation,
proves the formula for \(y_3\) in E.183.2.

The actual index-three mixed subgroup
\(\langle u,w^2\rangle\leq B_3\) has presentation
\[
             \langle u,z\mid(uz)^2=(zu)^2\rangle,
                    \qquad z=w^2.
\tag{E.183.7}
\]
This is the ordinary index-three Reidemeister--Schreier
presentation; its central cap is the group used in E.179.
The normalized mixed identity of E.179 is the equality
\(f(u,z^2)=f(u^2,z)\) in its full profinite completion.
By E.183.5 there is an actual homomorphism
\[
              \langle u,z\mid(uz)^2=(zu)^2\rangle
                      \longrightarrow B_4,
                       \qquad u\longmapsto c,\quad z\longmapsto\omega.
\tag{E.183.8}
\]
Its continuous extension therefore gives
\[
                         f(c,d)=f(c^2,\omega).
\tag{E.183.9}
\]
In particular this step needs no choice of a square root
of \(\omega\) in \(B_4\), and uses the mixed subgroup,
not a falsely marked copy of all of \(B_3\).

Let \(A_f\) be the standard genus-zero GT action on
\(\widehat B_4\), with the notation of
\citet[equation (1.3), p.657]{LitSchneps2006Automorphisms}.
At multiplier one its three generator formulas are
\[
 \begin{aligned}
 A_f(a)&=a,\\
 A_f(b)&=f(a^2,b^2)^{-1}b f(a^2,b^2),\\
 A_f(c)&=g^{-1}cg,
       \qquad g=f(y_3,c^2).
 \end{aligned}
\tag{E.183.10}
\]
The central boundary braid is fixed, so this action descends
to the cap in E.183.1.
Since \(a\) commutes with \(\omega,c\), both inputs of
\(g\) commute with \(a\).  Thus \(g\) commutes with \(a\).
Moreover the commuting factor \(a^{-2}\) in \(y_3\) may
be removed from an input of the commutator word \(f\):
form the homomorphism from the product of the two-generated
group with the central cyclic factor, and use the zero
exponent sums of \(f\).  Equations E.183.2 and E.183.9 give
\[
                 g=f(\omega,c^2)=f(c,d)^{-1}.
\tag{E.183.11}
\]
Consequently \(\operatorname{Ad}(g)A_f\) fixes both
\(a,c\), and its value on \(b\) is
\[
        g f(a^2,b^2)^{-1}b f(a^2,b^2)g^{-1}
                          =WbW^{-1}.
\tag{E.183.12}
\]
Comparison with E.183.4 proves equality of the two
automorphisms on the three generators, hence on the
entire completed group:
\[
                         \alpha_{f,P}=\operatorname{Ad}(g)A_f.
\tag{E.183.13}
\]
The element \(\omega\) and the square \(c^2\) are pure
braids, so \(g=f(\omega,c^2)\) is pure.  Thus the one
inner gauge preserves the actual branch ordering; no
independent geometric permutation or new GT parameter
has been introduced.  Restricting E.183.13 to the actual
open subgroup \(\widehat\Gamma_{0,5}\) proves the claim.

The program \path{code/check_gt_two_point_marking.py}
retains and verifies the displayed defining-relation path
and four related literal braid identities.  Its certificate
has five equalities and forty-three individual relation
steps, with SHA-256
\[
 \texttt{1cacc027efe676c6ec504960af38cf779}\,
 \texttt{1352c2bc046e9d43d3d00122733633c}.
\]
It is a word certificate in the presented braid group,
not a check in a finite braid image or a finite GT census.
\end{proof}
